La tercera y última técnica abandona las reglas explícitas y delega la generación lingüística a modelos de lenguaje grandes.

Este capítulo describe cómo se guía un LLM comercial mediante ingeniería de prompts~\cite{brown2020language} para resolver el mismo problema que los dos motores anteriores abordaron con gramáticas y embeddings.

\section{Diseño del prompt}
\label{sec:f3-prompt}

La ingeniería de prompts es el componente central de este motor. A diferencia de un modelo ajustado (\textit{fine-tuned})~\cite{brown2020language}, toda la adaptación al dominio CAA se codifica en el prompt enviado en cada llamada a la API.

\subsection{System prompt}

El \textit{system prompt} define el rol del modelo como motor de lenguaje de una aplicación CAA y establece once reglas estrictas que gobiernan la generación (véase Extracto de código~\ref{lst:system-prompt} del Anexo~\ref{anx:codigo}). Las reglas cubren cuatro aspectos.

La \textit{fidelidad semántica} (regla 1) es la restricción más importante: la oración debe expresar exactamente el significado de los pictogramas proporcionados, sin añadir información ni conceptos ausentes en la entrada. En un contexto CAA, generar contenido no seleccionado por el usuario equivaldría a poner en su boca palabras que no ha elegido.

Los \textit{elementos funcionales} (regla 2) especifican qué puede añadir el modelo: artículos, preposiciones de enlace, conjugación verbal y concordancia de género. Estas son las mismas operaciones que los motores anteriores resolvían con reglas explícitas (sección~\ref{sec:f1-postprocesamiento}); aquí se delegan al LLM mediante una instrucción textual.

El \textit{formato y registro} (reglas 3-5, 8, 10) restringe la salida a una única oración en registro coloquial-neutro, sin explicaciones ni alternativas. Sin estas restricciones, los LLMs frecuentemente producen texto adicional como ``Aquí tienes la oración:'' o generan múltiples variantes.

Las \textit{reglas lingüísticas} (reglas 6-7, 9, 11) cubren fenómenos específicos: interrogativas con signos ¿...?, posición de la negación, omisión del pronombre sujeto (\textit{pro-drop} del español) y concordancia de género. La regla 11, sobre concordancia, fue añadida respecto al diseño original para cubrir la categoría de concordancia del dataset de evaluación.

\subsection{Ejemplos few-shot}

Las reglas del \textit{system prompt} describen qué debe hacer el modelo, pero no le muestran cómo hacerlo. Sin ejemplos concretos, el modelo tiende a interpretar las instrucciones de forma demasiado literal o demasiado libre, produciendo salidas que cumplen las reglas individualmente pero no capturan el registro ni la brevedad esperados en CAA.

El motor utiliza 15 ejemplos \textit{few-shot}~\cite{brown2020language} que cubren las ocho categorías del dataset. Cada ejemplo es un par (pictogramas de entrada, oración esperada) que se inyecta como intercambio usuario-asistente en el historial de mensajes.

\begin{table}[H]
\centering
\begin{tabular}{r l l l}
\hline
\# & Categoría & Pictogramas & Oración \\
\hline
1 & Necesidad & quiero + agua & Quiero agua \\
2 & Necesidad & necesito + dormir & Necesito dormir \\
3 & Necesidad & quiero + comer + manzana & Quiero comer una manzana \\
4 & Inferencia & agua & Quiero agua \\
5 & Movimiento & ir + baño & Quiero ir al baño \\
6 & Estado & estoy + cansado & Estoy cansado \\
7 & Estado & tengo + hambre & Tengo hambre \\
8 & Descripción & perro + es + grande & El perro es grande \\
9 & Interrogativa & dónde + estar + mamá & ¿Dónde está mamá? \\
10 & Interrogativa & qué + comer & ¿Qué comes? \\
11 & Negación & no + quiero + dormir & No quiero dormir \\
12 & Negación & no + tengo + hambre & No tengo hambre \\
13 & Compuesta & agua + y + pan & Quiero agua y pan \\
14 & Compuesta & mañana + ir + colegio & Mañana quiero ir al colegio \\
15 & Concordancia & quiero + manzana + rojo & Quiero una manzana roja \\
\hline
\end{tabular}
\caption{Los 15 ejemplos \textit{few-shot} del motor de modelos de lenguaje.}
\label{tab:f3-fewshot}
\end{table}

El número de 15 se fijó porque es el mínimo que cubre las ocho categorías del dataset con al menos un ejemplo representativo; añadir más ejemplos incrementaría el coste de cada llamada sin beneficio proporcional, dado que los aproximadamente 1.204 tokens de entrada ya están dominados por el prompt fijo y los \textit{few-shot} existentes.

Los criterios de selección priorizan la cobertura de patrones que requieren inferencia: el ejemplo 4 enseña al modelo a expandir un solo sustantivo en una oración de necesidad; el ejemplo 5, a insertar la preposición contraída \textit{al}; el ejemplo 15, a corregir la concordancia de género. Los ejemplos se adaptaron al formato real del dataset, que utiliza tanto infinitivos como formas conjugadas como pictogramas de entrada.

\subsection{Construcción dinámica del prompt}

Cada mensaje de usuario formatea los pictogramas incluyendo su categoría gramatical entre paréntesis (por ejemplo, \texttt{Pictogramas: quiero (verbo modal) + agua (sustantivo)}).

Esta información desambiguadora procede del inventario de metadatos y ayuda al modelo a interpretar correctamente pictogramas polisémicos: ``como'' puede ser verbo (primera persona de ``comer'') o interrogativo (``¿cómo?''), y la categoría aclara la interpretación esperada.

Cada llamada consume aproximadamente 1.204 tokens de entrada, de los cuales la gran mayoría corresponden al \textit{system prompt} y los ejemplos \textit{few-shot}. La consulta del usuario aporta típicamente 10-20 tokens adicionales. Los tokens de salida son muy reducidos: 4 tokens en promedio, coherente con la brevedad de las oraciones de CAA.

\section{Metadatos de pictogramas}
\label{sec:f3-metadatos}

La temperatura de 0.3 prioriza la consistencia sobre la creatividad: en CAA, la variabilidad no aporta valor y la predecibilidad es deseable. Una temperatura de 0 sería completamente determinista, pero 0.3 permite cierta flexibilidad para manejar construcciones no vistas en los \textit{few-shot}.

El límite de 100 tokens es generoso para oraciones que promedian 4 tokens de salida, pero actúa como salvaguarda contra generaciones desbordadas. El timeout de 10 segundos refleja que, en una aplicación de comunicación asistida, el usuario espera activamente la respuesta.

El cliente implementa reintentos diferenciados según el tipo de error (véase Extracto de código~\ref{lst:llm-client} del Anexo~\ref{anx:codigo}): los errores de \textit{rate limit} se reintentan con backoff exponencial (2 y 4 segundos), los timeouts se reintentan inmediatamente, y los errores de servidor no se reintentan.

El coste se estima a partir de los precios de la API~\cite{openai2024pricing}: \$2.50/1M tokens de entrada y \$10.00/1M de salida para GPT-4o, frente a \$0.15/1M y \$0.60/1M para GPT-4o-mini.

Con los aproximadamente 1.204 tokens de entrada y 4 de salida descritos en la sección~\ref{sec:f3-prompt}, cada consulta cuesta \$0.00018 con GPT-4o-mini y \$0.00305 con GPT-4o, una diferencia de factor 16.7$\times$.

\section{Validación post-generación}
\label{sec:f3-validacion}

A diferencia de los motores basados en reglas, donde la corrección de la salida es consecuencia directa de las reglas aplicadas, un LLM puede producir respuestas vacías, omitir pictogramas de la entrada o ignorar instrucciones del prompt. En el contexto de CAA, donde el usuario depende del sistema para comunicarse, una salida que no refleje los pictogramas seleccionados es un fallo funcional. Por este motivo, el motor incorpora un validador específico que verifica dos aspectos que los LLMs no garantizan por diseño.

El primero es la cobertura de conceptos: el validador comprueba que cada pictograma seleccionado por el usuario esté representado en la oración generada, teniendo en cuenta las variaciones morfológicas del español (el pictograma ``querer'' puede aparecer como ``quiero'' o ``quiere'' en la salida). El segundo es el formato interrogativo: cuando la entrada contiene un pictograma interrogativo, se verifica que la respuesta incluya los signos de interrogación correspondientes, una convención ortográfica del español que el modelo no siempre respeta.

\section{Limitaciones}
\label{sec:f3-limitaciones}

El motor depende de la API de OpenAI, lo que implica conexión a Internet permanente y un servicio de pago que no es viable para uso offline. La latencia de red añade aproximadamente un segundo a cada generación, un tiempo perceptible en una aplicación de comunicación asistida.

Aunque el coste por consulta es reducido (\$0.00018 con GPT-4o-mini)~\cite{openai2024pricing}, el uso continuado genera un gasto acumulativo que no existe en los motores simbólicos.

Con temperatura 0.3, la misma entrada puede producir salidas ligeramente diferentes entre ejecuciones, lo que dificulta la depuración y la reproducibilidad de los resultados.

Las interrogativas son la categoría más difícil para el modelo, con las dos únicas interrogativas de los 15 ejemplos \textit{few-shot} insuficientes para cubrir la variedad de inferencias de persona verbal que esta categoría requiere.
